\section{Marco teórico}
\label{sec:marco-teorico}
El objetivo de seguridad de Olist Lab combina \textbf{confidencialidad}, \textbf{integridad} y \textbf{disponibilidad}. La confidencialidad limita quién puede leer datos y registros. La integridad permite detectar cambios no autorizados en una copia. La disponibilidad requiere que el equipo pueda recuperar la base tras una pérdida. Usamos NIST CSF 2.0 como referencia para relacionar estas metas con identificación de riesgos, protección, detección, respuesta y recuperación \cite{nistcsf}.

\paragraph{Cifrado y autenticación de copias.} SQLCipher cifra las páginas de la base con AES-256 y protege su integridad con HMAC según su configuración por defecto \cite{sqlcipher}. Esto protege el archivo frente a lectura directa sin la clave. Para respaldos, AES-GCM cifra y autentica la copia. La etiqueta permite detectar alteraciones y el nonce debe ser único para cada cifrado con una misma clave \cite{gcm}. La prueba del final muestra que la aplicación admite una copia válida y rechaza una alterada.

\paragraph{Hash de contraseñas y TLS.} Un hash de contraseñas con sal y costo de cómputo permite verificar una contraseña sin guardarla de forma recuperable \cite{owasp}. La sal dificulta reutilizar cálculos entre cuentas. TLS protege el tráfico en tránsito y puede autenticar al servidor si se valida su certificado \cite{tls}. En la entrega final el certificado del servidor lo emite una CA propia del proyecto (RFC 5280 \cite{x509}); el cliente que instala esa raíz valida la cadena y el nombre del host, mientras que uno sin ella rechaza la conexión.

\paragraph{Autenticación, autorización y trazabilidad.} La autenticación establece qué cuenta inició sesión. La autorización comprueba qué puede hacer esa cuenta. En Olist Lab, Flask permite a lectores consultar los indicadores y reserva la auditoría a administradores. Los eventos permiten reconstruir accesos, consultas y fallos. Su valor depende de limitar el acceso y acordar plazos de conservación.

\paragraph{Minimización y reidentificación.} Quitamos los identificadores del cliente de las salidas y agrupamos los resultados. La vinculación por fechas, ubicación o identificadores persistentes puede permitir reidentificación si alguien dispone de otra fuente. Una agrupación pequeña también puede revelar información por inferencia. Por eso la DPIA evalúa el conjunto original, la base y la salida pública por separado.

\section{Requerimientos de seguridad}
\label{sec:requisitos-seguridad}
La tabla resume cada requerimiento de seguridad, en qué entrega se implementó (E1: avance; EF: entrega final) y la evidencia reproducible. Las comprobaciones marcadas con \texttt{pytest} se ejecutan sobre una instancia aislada con datos sintéticos; las marcadas con \texttt{verify\_live} se ejecutan sobre la instancia real mediante HTTPS validado contra la CA propia (sección~\ref{sec:validacion}).

{\small
\begin{longtable}{@{}P{.055\textwidth}P{.26\textwidth}P{.06\textwidth}P{.29\textwidth}P{.23\textwidth}@{}}
\caption{Requerimientos de seguridad, entrega, evidencia y alcance.}\\
\toprule
\textbf{ID} & \textbf{Requerimiento y control} & \textbf{Ent.} & \textbf{Evidencia observada} & \textbf{Alcance o límite}\\
\midrule
\endfirsthead
\toprule
\textbf{ID} & \textbf{Requerimiento y control} & \textbf{Ent.} & \textbf{Evidencia observada} & \textbf{Alcance o límite}\\
\midrule
\endhead
S01 & Datos en reposo: base SQLCipher (AES-256 + HMAC por página); ZIP fuente sellado con AES-256-GCM; claves de 256 bits con permisos 0600 y carpeta configurable (\texttt{OLIST\_KEY\_DIR}). & E1 / EF & Cabecera no legible; una clave errónea no abre la base; ZIP en claro eliminado; 96\,470 filas con la clave. & Las claves siguen en el equipo de la demo salvo que se use otro volumen.\\
S02 & Datos en tránsito: HTTPS con certificado emitido por la CA propia del proyecto, TLS $\geq$ 1.2 y HSTS. & EF & TLS 1.3 (\texttt{TLS\_AES\_256\_GCM\_SHA384}) verificado contra la CA; un cliente sin la CA rechaza la conexión. & La CA es privada: solo los equipos que la instalen confían en ella.\\
S03 & Hash de contraseñas: scrypt con sal (Werkzeug); mínimo 12 caracteres con mayúsculas, minúsculas y dígito. & E1 / EF & Hashes únicos por cuenta; contraseñas débiles rechazadas. & Sin MFA.\\
S04 & Gestión de accesos: sesión obligatoria, roles \texttt{admin}, \texttt{analyst} y \texttt{reader}, CSRF, sesión de 30 min, bloqueo por usuario (5) y por IP (20) en 15 min, cookies \texttt{Secure}, \texttt{HttpOnly} y \texttt{SameSite}, CSP. & E1 / EF & API sin sesión 401; panel 302; auditoría 403/200; sin CSRF 400; sexto intento 429. & Los roles \texttt{analyst} y \texttt{reader} comparten permisos analíticos.\\
S05 & Registros de auditoría: ingresos, fallos, bloqueos, consultas, filtros inválidos, accesos denegados, CSRF rechazado y altas/bajas de usuarios; retención de 180 días. & E1 / EF & Eventos persistidos sin contraseñas; \texttt{purge-audit} elimina registros vencidos. & Los logs están en la misma base; no hay SIEM.\\
S06 & Respaldo y recuperación: copia AES-256-GCM autenticada, rotación (\texttt{-{}-keep}), restauración con \texttt{integrity\_check}. & E1 / EF & 96\,470 filas restauradas; una copia con un bit alterado se rechaza sin escribir salida. & La copia debe guardarse fuera del equipo.\\
S07 & Herramientas de análisis: Bandit (SAST), pip-audit y npm audit (SCA), gitleaks (secretos), pytest y verificación en vivo; CI semanal. & EF & 0 hallazgos tras corregir 4 B101, 3 B608 y 15 CVE de npm (tabla~\ref{tab:herramientas}). & No sustituye una prueba de penetración.\\
S08 & Uso seguro: política, procedimiento, concientización y formación del equipo. & E1 / EF & Sección~\ref{sec:politicas}; ejercicio de formación ejecutado sobre el incidente real de la sección~\ref{sec:incidente-real}. & El cumplimiento organizativo no se audita externamente.\\
S09 & Riesgos de una brecha y entidades afectadas. & E1 / EF & Modelo de amenazas y DPIA (secciones~\ref{sec:amenazas} y~\ref{sec:dpia}). & Evaluación cualitativa.\\
S10 & Plan de respuesta ante incidentes y recomendaciones futuras. & E1 / EF & Plan NIST SP 800-61 aplicado a un caso real; recomendaciones en la sección~\ref{sec:mejoras}. & Plazos legales no evaluados por un abogado.\\
\bottomrule
\end{longtable}
}
